% generator_I14.tex -- Prop I.14 (Theor.): two right angles on one side => straight line.
\documentclass[12pt]{article}\usepackage{geometry}\geometry{a5paper,margin=1.5cm}
\usepackage[lining]{ebgaramond}\usepackage{amssymb}\usepackage{byrne}
\def\mpPre{u := 1cm; textLabels := false;}
\begin{document}
\defineNewPicture[1/4]{%
  pair A,B,C,D,E;%
  A:=(u,5/2u);B:=(-7/4u,0);C:=B xscaled -1;D:=(0,0);%
  E:=(xpart(C),-1/2ypart(A));%
  byAngleDefine(B,D,A,byyellow,SOLID_SECTOR);%
  byAngleDefine(C,D,A,byblue,SOLID_SECTOR);%
  byAngleDefine(E,D,C,byred,SOLID_SECTOR);%
  draw byNamedAngleResized();%
  draw byLine(A,D,byred,SOLID_LINE,REGULAR_WIDTH);%
  draw byLine(E,D,byyellow,SOLID_LINE,REGULAR_WIDTH);%
  draw byLine(B,D,byblue,SOLID_LINE,REGULAR_WIDTH);%
  draw byLine(C,D,byblack,SOLID_LINE,REGULAR_WIDTH);%
  draw byLabelsOnPolygon(E,D,B,noPoint)(ALL_LABELS,0);%
  draw byLabelsOnPolygon(C,B,noPoint)(OMIT_LAST_LABEL,0);%
  draw byLabelLineEnd(A,D,0);%
}
\noindent\begin{minipage}[t]{0.45\linewidth}\centering\vspace{0pt}%
\drawCurrentPicture\end{minipage}\hfill
\begin{minipage}[t]{0.52\linewidth}\vspace{0pt}%
\textbf{Book~I.\enspace Prop.\ XIV. Theor.}
\medskip\noindent\itshape
If two straight lines (\drawUnitLine{BD} and \drawUnitLine{DC}),
meeting a third straight line (\drawUnitLine{AD}), at the same point,
and at opposite sides of it, make with it adjacent angles
(\drawAngle{BDA} and \drawAngle{CDA}) equal to two right angles;
these straight lines lie in one continuous straight line.
\bigskip\noindent\upshape
For, if possible let \drawUnitLine{ED}, and not \drawUnitLine{DC},\\
be the continuation of \drawUnitLine{BD},\\
then $\drawAngle{BDA}+\drawAngle{CDA,EDC}=\drawTwoRightAngles$\\[4pt]
but by the hypothesis $\drawAngle{BDA}+\drawAngle{CDA}=\drawTwoRightAngles$\\[4pt]
$\therefore\drawAngle{CDA,EDC}=\drawAngle{CDA}$~(ax.~III); which is absurd~(ax.~IX).\\[4pt]
$\therefore \drawUnitLine{ED}$ is not the continuation of \drawUnitLine{BD},
and the like may be demonstrated of any other straight line except
\drawUnitLine{DC}, $\therefore \drawUnitLine{DC}$ is the continuation of
\drawUnitLine{BD}.
\begin{flushright}Q.\ E.\ D.\end{flushright}
\end{minipage}
\end{document}
